\documentclass[12pt,a4paper]{article}
\usepackage[utf8]{inputenc}
\usepackage{xepersian}
\usepackage{graphicx}
\usepackage{hyperref}

\settextfont{XB Niloofar}
\setdigitfont{XB Niloofar}

\title{فهرست مطالب}
\author{}
\date{}

\begin{document}

\maketitle

\textbf{مصاحبهٔ طراحی سیستم: راهنمای یک فرد داخلی}

\textit{پیشگفتار}

\begin{itemize}
\item \textbf{فصل ۱:} مقیاس‌پذیری از صفر تا میلیون‌ها کاربر
\item \textbf{فصل ۲:} تخمین‌های سرانگشتی
\item \textbf{فصل ۳:} چارچوبی برای مصاحبه‌های طراحی سیستم
\item \textbf{فصل ۴:} طراحی یک محدودکنندهٔ نرخ\footnote{Rate Limiter}
\item \textbf{فصل ۵:} طراحی هَش سازگار\footnote{Consistent Hashing}
\item \textbf{فصل ۶:} طراحی یک ذخیره‌ساز کلید-مقدار\footnote{Key-Value Store}
\item \textbf{فصل ۷:} طراحی یک تولیدکنندهٔ شناسهٔ منحصربه‌فرد در سیستم‌های توزیع‌شده
\item \textbf{فصل ۸:} طراحی یک کوتاه‌کنندهٔ لینک\footnote{URL Shortener}
\item \textbf{فصل ۹:} طراحی یک خزندهٔ وب\footnote{Web Crawler}
\item \textbf{فصل ۱۰:} طراحی یک سیستم اطلاع‌رسانی\footnote{Notification System}
\item \textbf{فصل ۱۱:} طراحی یک سیستم فید خبری\footnote{News Feed System}
\item \textbf{فصل ۱۲:} طراحی یک سیستم چت
\item \textbf{فصل ۱۳:} طراحی یک سیستم تکمیل خودکار جستجو\footnote{Search Autocomplete}
\item \textbf{فصل ۱۴:} طراحی یوتیوب
\item \textbf{فصل ۱۵:} طراحی گوگل درایو
\item \textbf{فصل ۱۶:} یادگیری ادامه دارد
\end{itemize}

\textit{سخن پایانی}

\section*{توضیحات}

\begin{itemize}
\item اصطلاحات تخصصی در فوت‌نوت‌ها به انگلیسی ذکر شده‌اند.
\item از واژه‌های رایج در صنعت فناوری اطلاعات استفاده شده است.
\item ساختار فهرست مطالب اصلی حفظ شده است.
\end{itemize}

\end{document}
